\documentclass[10pt]{article}
\usepackage[letterpaper,margin=0.72in]{geometry}
\usepackage{booktabs,longtable,array,tabularx}
\usepackage{xcolor}
\usepackage{hyperref}
\usepackage{microtype}
\usepackage{enumitem}
\usepackage{amsmath,amssymb}
\definecolor{navy}{HTML}{153A5B}
\definecolor{lightblue}{HTML}{EAF3F8}
\definecolor{darkgray}{HTML}{444444}
\hypersetup{colorlinks=true,linkcolor=navy,urlcolor=navy}
\setlength{\parindent}{0pt}
\setlength{\parskip}{5pt}
\setlist{nosep,leftmargin=1.3em}
\renewcommand{\arraystretch}{1.12}
\newcommand{\obs}{\textbf{Observed.}\ }
\newcommand{\inferlabel}{\textbf{Inferred.}\ }
\newcommand{\hyp}{\textbf{Hypothesis.}\ }
\newcommand{\na}{\textemdash}

\title{\textbf{CAROM Frozen-GPT-2 Compiled-Channel Experiment}\\
Detailed protocol, checkpoint results, causal interventions, and next-run design}
\author{CAROM research record}
\date{21 July 2026}

\begin{document}
\maketitle

\begin{center}
\fcolorbox{navy}{lightblue}{\parbox{0.91\linewidth}{
\textbf{Bottom line.} A single frozen-GPT-2-small CAROM arm learned a moderately
accurate in-distribution executor (0.412 slot accuracy at 4,000 updates) with
strongly ordered-looking trajectories ($\tau=0.904$). The natural trajectory
became causally useful: replacing it by a shuffled schedule reduced accuracy
from 0.412 to 0.297. This is meaningful mechanistic evidence, but accuracy is
well below the matched TinyLM arm (0.709), and held-out length-5 transfer remains
weak (0.214; 0.234 with a longer inference budget). A longer frozen-GPT-2 run is
therefore justified as an optimization test, not yet as evidence of structural
generalization or learned commutation.}}
\end{center}

\section{Question and scope}

The experiment asked whether natural-language command spans encoded by a
\emph{frozen} GPT-2-small model can drive a CAROM compiled heteroclinic channel.
The task is synthetic compositional sequence transformation. Each example has
six workspace slots containing symbols in $\mathbb{Z}_8$ and a program of two
to five primitive transforms. Command descriptions are presented in scrambled
surface order; the model must infer dependencies, produce a command-activity
trajectory, mix reusable neural operators, and emit the final six symbols.

Training used program lengths 2--4. Length 5 was held out structurally. Thus:

\begin{itemize}
\item L2--4 accuracy measures in-distribution compositional execution;
\item L5 accuracy measures extrapolation to a longer program;
\item edge accuracy measures dependency-graph prediction;
\item itinerary diagnostics describe the ordering of dominant command modes;
\item trajectory replacements test whether the learned activity trajectory is
causally relevant rather than merely correlated with the answer.
\end{itemize}

This was one seed and one frozen evaluation corpus. It is a mechanism screen,
not a population estimate.

\section{Architecture and training protocol}

\subsection{Language-conditioned compiler}

GPT-2-small (124M parameters, hidden width 768, final hidden layer) was frozen.
Span representations were produced using offset-mapped token spans. CAROM then
used learned heads to (i) predict a directed dependency graph between command
spans, (ii) compile that graph into generalized Lotka--Volterra inhibition, and
(iii) route command representations onto a shared bank of neural workspace
operators. Only the CAROM-side parameters were trained.

The workspace has six slots, vocabulary size eight, and up to five commands.
The recurrent controller executes 72 internal steps at training/default
inference. The endpoint objective is slotwise cross-entropy. The training loss
also includes directed-edge binary cross-entropy and an entry-command term.

\subsection{Exact run configuration}

\begin{tabularx}{\linewidth}{>{\bfseries}l X}
\toprule
Seed & 0 \\
Training updates & 4,000 \\
Batch size & 64 \\
Optimizer & AdamW, learning rate $2\times10^{-3}$, weight decay $10^{-4}$ \\
Schedule & OneCycleLR over 4,000 updates \\
Gradient clipping & Global norm 1.0 \\
Edge-loss coefficient & 1.0; entry term coefficient 0.2 \\
Checkpoints & 0, 500, 1,000, \ldots, 4,000 \\
Frozen L2--4 corpus & 256 examples; corpus seed 1234 \\
Frozen L5 corpus & 128 examples; corpus seed 5678 \\
Hardware & RunPod Secure Cloud, one NVIDIA A100 SXM4 80GB \\
Image & \texttt{runpod/pytorch:1.0.2-cu1281-torch280-ubuntu2404} \\
Price and lifetime & USD 1.49/hr; approximately 2h09m including orchestration and repair \\
\bottomrule
\end{tabularx}

Training itself reached step 4,000 in 1,897 seconds (31.6 minutes). The original
post-training pass then failed on malformed Python corpus unpacking. The repair
loaded the missing \texttt{SP}, \texttt{E}, and \texttt{ENT} fields explicitly
and added an \texttt{--interventions-only} continuation. It reused all retained
checkpoints, evaluated nine checkpoints in about 40 seconds, and exited zero.

\section{Evaluation instruments}

\paragraph{Repaired itinerary panel.}
At each internal step, the dominant live command mode is recorded; consecutive
duplicates are collapsed into phases. Coverage is the fraction of true command
modes visited. Exact order requires the deduplicated visited sequence to equal
the true topological order. Kendall-style $\tau$ scores concordant versus
discordant rank pairs among visited modes. Because incomplete trajectories can
still have high $\tau$, tau must be read beside coverage and exact order.

\paragraph{Trajectory interventions.}
The natural learned trajectory is compared with three replacements on exactly
the same frozen L2--4 examples:

\begin{itemize}
\item \textbf{forced}: true command order, equal dwell segments, amplitude 1.5;
\item \textbf{shuffled}: an explicitly non-true command order with equal dwell;
\item \textbf{smeared}: all live commands active simultaneously with equal mass.
\end{itemize}

Natural-minus-shuffled accuracy is a causal sensitivity diagnostic. The forced
condition is not a perfect oracle: it changes amplitude and dwell statistics as
well as order, so forced below natural is interpretable rather than contradictory.

\paragraph{Integration-budget sweep.}
The same L5 examples are evaluated with 72, 100, and 120 internal steps. A large
increase would implicate insufficient settling time; a small increase leaves a
structural or learned-computation limitation as the main explanation.

\section{Checkpoint results}

\begin{longtable}{rrrrrrrr}
\caption{Frozen GPT-2 checkpoint panel. Accuracy is mean slot accuracy.}\\
\toprule
Step & L2--4 acc & Edge acc & $\tau$ & Coverage & Exact & L5 acc & Nat$-$shuf \\
\midrule
\endfirsthead
\toprule
Step & L2--4 acc & Edge acc & $\tau$ & Coverage & Exact & L5 acc & Nat$-$shuf \\
\midrule
\endhead
0    & .139 & .297 & .040 & .383 & .000 & .135 & .000 \\
500  & .206 & .703 & .477 & .500 & .043 & .163 & -.003 \\
1000 & .240 & .703 & .945 & .695 & .309 & .176 & .003 \\
1500 & .266 & .703 & .905 & .682 & .281 & .145 & -.007 \\
2000 & .291 & .688 & .862 & .663 & .246 & .217 & .025 \\
2500 & .331 & .702 & .871 & .662 & .262 & .232 & .060 \\
3000 & .372 & .705 & .892 & .678 & .273 & .221 & .079 \\
3500 & .402 & .709 & .904 & .700 & .285 & .225 & .089 \\
4000 & .412 & .715 & .904 & .710 & .297 & .214 & .115 \\
\bottomrule
\end{longtable}

\obs Accuracy rose from 0.139 to 0.412. Edge accuracy jumped to about 0.70 by
step 500 and then largely plateaued; endpoint accuracy continued improving.
Itinerary tau rose sharply and remained near 0.86--0.90 after step 1,500, while
coverage slowly recovered to 0.710. Exact full order reached only 0.297.

The high tau is therefore not equivalent to complete itinerary execution: the
model orders the modes it visits fairly well but omits or fails to dominate with
some modes. This distinction explains how tau can look strong while endpoint
accuracy remains moderate.

\section{Causal trajectory results}

\begin{longtable}{rrrrrr}
\caption{Endpoint accuracy under natural and replaced trajectories.}\\
\toprule
Step & Natural & Forced & Shuffled & Smeared & Natural$-$shuffled \\
\midrule
\endfirsthead
\toprule
Step & Natural & Forced & Shuffled & Smeared & Natural$-$shuffled \\
\midrule
\endhead
0    & .139 & .139 & .139 & .139 & .000 \\
500  & .206 & .209 & .209 & .208 & -.003 \\
1000 & .240 & .236 & .238 & .236 & .003 \\
1500 & .266 & .264 & .273 & .263 & -.007 \\
2000 & .291 & .272 & .266 & .285 & .025 \\
2500 & .331 & .278 & .272 & .264 & .060 \\
3000 & .372 & .311 & .292 & .308 & .079 \\
3500 & .402 & .317 & .314 & .318 & .089 \\
4000 & .412 & .309 & .297 & .301 & .115 \\
\bottomrule
\end{longtable}

\obs The intervention gap is essentially zero early and grows after step 2,000.
At step 4,000, shuffling lowers accuracy by 11.5 percentage points relative to
the natural trajectory. Smeared and forced schedules are similarly below
natural. The learned trajectory has therefore become load-bearing for endpoint
performance on the measured corpus.

\inferlabel This rules out the strongest ``decorative trajectory'' account at the final
checkpoint. It does not prove that every visited phase or every ordering detail
is necessary, nor that the mechanism extrapolates beyond this task distribution.

\section{Held-out length and integration depth}

\begin{longtable}{rrrrr}
\caption{L5 endpoint accuracy versus inference steps.}\\
\toprule
Step & Default L5 & $S=72$ & $S=100$ & $S=120$ \\
\midrule
\endfirsthead
\toprule
Step & Default L5 & $S=72$ & $S=100$ & $S=120$ \\
\midrule
\endhead
0    & .135 & .135 & .135 & .135 \\
500  & .163 & .163 & .164 & .163 \\
1000 & .176 & .176 & .180 & .182 \\
1500 & .145 & .145 & .155 & .168 \\
2000 & .217 & .217 & .233 & .243 \\
2500 & .232 & .232 & .237 & .249 \\
3000 & .221 & .221 & .224 & .230 \\
3500 & .225 & .225 & .227 & .224 \\
4000 & .214 & .214 & .224 & .234 \\
\bottomrule
\end{longtable}

\obs At the final checkpoint, adding 48 internal steps improves L5 accuracy by
2.1 points. The best measured L5 value (0.249) occurs at step 2,500 with
$S=120$, still far below in-distribution performance.

\inferlabel Insufficient integration depth contributes slightly, but it is not the main
L5 failure. The dominant limitation is learned computation/representation or
structural extrapolation, not simply a short settling horizon.

\section{Comparison with the retained TinyLM arm}

The TinyLM comparison used the same nine checkpoint indices and intervention
definitions, but its language encoder, batch/optimization details, and training
artifact are not identical. It is contextual, not a controlled claim that one
language model is intrinsically superior.

\begin{tabular}{lrr}
\toprule
Final metric & TinyLM & Frozen GPT-2 \\
\midrule
L2--4 accuracy & .709 & .412 \\
Itinerary $\tau$ & .825 & .904 \\
Coverage & .911 & .710 \\
Natural$-$shuffled & .227 & .115 \\
L5 accuracy & .356 & .214 \\
L5 at $S=120$ & .362 & .234 \\
\bottomrule
\end{tabular}

GPT-2 shows higher ordering correlation among visited phases but markedly lower
coverage, accuracy, causal sensitivity, and length transfer. A plausible
optimization explanation is that the CAROM routing heads must extract useful
signals from a wider, frozen 768-dimensional space under the same 4,000-update
budget. This is a hypothesis, not established by the current arm; a longer run
is the direct test.

Importantly, these numbers do \emph{not} establish learned commutation. Learned
commutation predicts declining order sensitivity, which requires the repaired
nonlinear swap-discrepancy and JVP commutator-energy probes. Here the observed
natural-minus-shuffled gap grows, so the measured trajectory becomes more, not
less, causally consequential over training.

\section{What the experiment establishes}

\begin{itemize}
\item \obs A frozen GPT-2 span encoder can support learning above the random
endpoint baseline in this CAROM compiled-channel architecture.
\item \obs Dependency prediction is learned early, but endpoint execution
continues improving after edge accuracy plateaus.
\item \obs The final natural trajectory is causally useful relative to shuffled,
smeared, and fixed-dwell forced replacements.
\item \obs High tau coexists with incomplete coverage and low exact-order rate;
tau alone would overstate mechanism quality.
\item \obs Extra inference depth only modestly improves held-out L5 accuracy.
\item \textbf{Not established:} robust GPT-2 semantic compilation, multi-seed
reliability, learned commutation, schedule-uniform stability, or structural
length generalization.
\end{itemize}

\section{Higher-accuracy follow-up}

The cleanest next test changes only the training budget: frozen GPT-2-small,
seed 0, batch 64, losses and architecture unchanged, but 12,000 updates with a
OneCycle schedule defined over all 12,000 updates and checkpoints every 1,000.
This triples the number of training examples from 256,000 to 768,000 while
preserving comparability. The same frozen corpora and intervention suite will be
used. Primary success is final L2--4 accuracy above 0.55; a stronger target is
0.65. Mechanism gates remain coverage, exact order, natural-minus-shuffled, and
L5. If accuracy rises while these degrade, the optimization succeeds only at
the endpoint and the mechanism claim weakens.

Based on the measured 31.6-minute training time for 4,000 steps, expected
training is about 95 minutes plus setup/checkpoint evaluation. The proposed
guardrail is one Secure Cloud A100 SXM4 80GB at USD 1.49/hr, expected two hours
(USD 2.98), hard cap three hours (USD 4.47), with auto-termination, immediate
artifact retrieval and hash verification, then termination.

\section{Provenance}

Primary results are
\nolinkurl{projects/carom/experiments/20260721T222400Z-carom-gpt2-arm/}.
The retrieved package contains nine checkpoints, nine per-checkpoint JSON files,
\texttt{summary.json}, logs, repaired source, and a SHA-256 manifest covering 20
artifact files (105 MiB). The repaired runner SHA-256 is
\begin{quote}\small\ttfamily
211ae1387deec50c70d5f60b3571f17c\\
8e89d8ebc555763ef78d1be5d2ac8cc0
\end{quote}
The compiled-channel source SHA-256 is
\begin{quote}\small\ttfamily
9171bcd065226d3f5ca96d7bb8f9f20d\\
9f144a11f251923382cbefb7be112223
\end{quote}

\end{document}
